% Separate audit fragment; requires amsmath, amssymb and hyperref.
% No unconditional weak-base PA-to-calculus theorem is asserted here.
\section{Focused weak-base calibration audit}
\label{sec:focused-calibration-audit}

The inspected source is the repository at commit
\path{1c9524d8bf2e0a27b9a16ff622359927f536e663}.
The original written calibration, in
\path{report/reference/original_pa195_calibration.tex} of the
\path{PA/docs/pa194_conpa_2026-11-13} package, proves its endpoint over
PA. It asserts primitive-recursive proof transformations with
PA-provable correctness. This does not alone yield correctness in
intuitionistic primitive recursive arithmetic.

There is nevertheless a concrete improvement. The bounded all-store
proof in the present command-preservation argument replaces the role of
the full-PA induction in the upstream
\path{PA196/Arithmetic/CompilerInvariant.lean}, theorem
\path{fuel_invariant}. It applies to arbitrary internal finite
numerical derivations and the new commands. The remaining calibration
question is earlier: can every internal PA refutation be converted into
such a numerical derivation over the weaker base?

\subsection{A finite-demand induction principle}

The following sufficient principle makes the uniformity requirement
precise. Its hypotheses must be instantiated; the existence of such a
principle does not automatically verify a compiler.

Code a request by a pair $(r,x)$, with numerical rank $r$ and parameter
code $x$. Let $\operatorname{Req}(r,x)$ and
$\operatorname{Cert}(r,x,c)$ be primitive-recursive Boolean predicates.
Suppose a primitive-recursive function $Q(r,x)$ returns the finite list
of child requests. For a valid parent, each child is valid and has rank
strictly below $r$. Crucially, $Q$ is computed from the request before
receiving any child certificates. Let a primitive-recursive constructor
$K$ combine certificates for those children. Require that iPRA proves
the just-stated rank and validity properties and proves
\[
 \operatorname{Req}(r,x)\ \wedge\
 \bigwedge_{(j,y)\in Q(r,x)}
       \operatorname{Cert}(j,y,c_{j,y})
 \ \longrightarrow\
 \operatorname{Cert}\bigl(r,x,K(r,x,(c_{j,y}))\bigr).
\]
The empty-list case includes the actual uniform leaf suppliers.

Fix a root rank $D$ and parameter bound $s$. Define
\[
 F_D(v)=1+\max\left(\{v\}\cup
  \{y: r\le D,\ x\le v,\ (j,y)\in Q(r,x)\}\right),
 \qquad B(r)=F_D^{\,D-r}(s)\quad(r\le D).
\]
These are primitive-recursive functions: the maximum is a bounded scan
over finite computed lists. If $x\le B(r)$ and $(j,y)\in Q(r,x)$
with $j<r$, then
\[
 y\le F_D(B(r))\le B(j).
\]
Construct a table in ascending rank, containing a certificate for
every $x\le B(r)$ in row $r$. Use the already completed rows for child
lookups; invalid requests receive a default output. The constructor is
primitive recursive, regardless of how large its certificate outputs
are. The invariant after completing rows below $k$ is
\[
 \bigwedge_{r<k}\ \bigwedge_{x\le B(r)}
 \left(\operatorname{Req}(r,x)\to
       \operatorname{Cert}(r,x,\operatorname{table}[r,x])\right).
\]
Both conjunctions are finite Boolean folds, so this is a
quantifier-free induction invariant. An inner bounded Boolean
induction constructs the next row. The child-bound inequality and the
local correctness hypothesis prove that row's validity. Hence iPRA
proves the root certificate correct, uniformly in $D,s$.

Thus the unbounded oracle formulations used in the existing proofs
need not be retained. But the principle requires the complete finite
request list and the actual local suppliers. A statement that all
recursive calls have ``some sufficiently large bound'' does not
instantiate these hypotheses.

\subsection{Concrete request maps and the remaining supplier}

The following source paths are relative to
\path{PA/jamesyc-upstream/pa196/PA196/Arithmetic/Transport/}.
The endpoint \path{ordinary_derivation_requested} in
\path{NumericPAeqSequentDerivation.lean} uses the following interface.
Write $g=\operatorname{guardedCode}(f,S)$ and
$P(n,l,z)=\operatorname{parameters}(0,n,l,z)$. Then
\[
\begin{aligned}
 \operatorname{RequestedCertificate}(f,S):={}&
 \forall n,l,z\,\bigl(\operatorname{Covered}(0,n,l,z,g)\\
 &\to \exists c\,
  \operatorname{verify}(c)=\operatorname{translate}(P(n,l,z),g)+1\bigr).
\end{aligned}
\]
The existing full induction quantifies all environments and smaller
derivations. Several finite request maps can be read directly from its
local rules. In the existential rule, the padding budget is
\[
 r=\max\bigl(\operatorname{guardedCode}
      (f,\operatorname{insert}(\operatorname{substs1}(u,p),S)),
      \operatorname{existentialRuleCode}(f,p,u)\bigr),
\]
and the child is queried at
$(\operatorname{padEnvironment}(n,r),l,z+l+r)$.
The shift rule has the same map with
$r=\operatorname{guardedCode}(f,S)$ at the child environment.
The theorem \path{requested_universal_certificate} in
\path{FoundationRelationalSequentRequestedUniversal.lean} uses
$d=\operatorname{free}(p)+1$ and queries its opened child at
\[
 \bigl(l::\operatorname{padEnvironment}(n,d),\ l+1,\ z+l+d\bigr).
\]
These maps are independent of returned certificate values. They are
concrete ingredients for the preceding lemma. Nested composition and
universal-context helpers induce additional requests whose coverage
must also be proved.

The axiom branch calls
\path{induction_axiom_requested_certificate} in
\path{NumericPAeqInductionAxiomCertificate.lean}. For every internally
coded PA induction axiom $q$, its target is a certificate for
\[
 \operatorname{translate}\bigl(\operatorname{parameters}(0,n,l,z),
       \operatorname{formulaSyntax}(\operatorname{parameters}(0,f),q)
                         \bigr)
\]
at every covered triple $(n,l,z)$. This closed-axiom theorem permits
arbitrary $f$; a closed-environment lemma removes its influence.
Its current proof assumes a model of full PA. Unfolding it reveals
further explicit full-induction uses:

\begin{itemize}
\item \path{NumericPAeqGuardedAlls.lean},
\path{matrixSupply_alls}, uses
\path{InductionScheme.succ_induction} with
\texttt{C := Set.univ}. Its invariant quantifies all matrices and
environments and propagates a certificate-supplier interface through
arbitrarily many universal closures.
\item \path{NumericPAeqCanonicalMatrixCertificate.lean},
\path{matrix_certificate}, invokes
\path{open_source_formula_realization_certificate}.
In \path{NumericPAeqFormulaRealization.lean},
\path{open_source_formula} proves this by another
\texttt{C := Set.univ} induction over formula codes. Its
\texttt{FormulaSupply} quantifies the changing depths, substitution
vectors, environments, and requested labels.
\item \path{RelationalPredicateTemplateOrdinalRequested.lean}
uses the fixed derivation \texttt{inductionDerivation} through
\path{template_nd_requested_certificate}, with
\path{induction_axiom_certificate} as callback. The fixed template
is promising, but uniform substitution, coverage, and callback
correctness are still required.
\end{itemize}

The two full-induction sites highlighted in the forward-audit README
are consequently examples, not an exhaustive transitive list.
Removing the top sequent induction and the command-fuel induction
alone does not close the forward calibration.

\subsection{Exact conclusion of the focused attempt}

No obstruction or counterexample to weak-base calibration was found.
In particular, a generated PA proof may contain an arbitrary PA
induction instance while its syntactic producer is verified in a much
weaker metatheory; the induction strength of the output proof must not
be confused with that of its verifier.

However, a complete iPRA or $I\Sigma_1$ proof of the PA-to-numerical
calibration is not supplied here. The precise remaining construction
is the uniform arbitrary-matrix supplier above, with a complete
finite-demand and coverage proof for its formula-realization,
closure, and sequent dependencies. The bounded-demand lemma explains
how suitable explicit suppliers would remove the high-induction
formulations. It cannot replace those suppliers. The numerical
verifier and replacement-interface maps must additionally be matched
to the exact new calculus. The physical simulation is a separate
proof, supplied in Section~\ref{sec:physical-ipra}, and does not
complete this compiler. Therefore the new abstract-command theorem does
not by itself justify the unconditional machine-specific iPRA
equivalence requested by the user.
